\section{Method}
\label{sec:method}

\subsection{Re-deriving the eligible frame}
\label{sec:method:frame}

The study inherited a corpus of four peer-reviewed studies. That corpus is a property of the
search that produced it rather than a census of the field, and its own record says so: of
eighteen search channels, seven returned HTTP 403, 418 or 429 before a result page,
including both venue databases where software-engineering reproduction work appears and both
citation traversals. The screening log records no excluded study, so no candidate pool was
ever narrowed to four; the four entered by artifact-link traversal from seed benchmarks.

We re-ran discovery over channels that answer -- OpenAlex, Crossref, DBLP, arXiv and
OpenReview -- with six queries that do not require a paper to name SWE-bench. OpenAlex and
Crossref index the ACM and IEEE proceedings whose own endpoints still refuse an automated
client. Every attempt records the channel's reported total against what was actually
retrieved, so a partial tail is visible rather than passing as complete.

Discovery yielded 27{,}921 deduplicated records. A deterministic first-stage gate removed
records outside the publication window or carrying no topical signal, and a second layer
required a term specific to repository-level software maintenance rather than the generic
\emph{benchmark} or \emph{evaluation}. That layer is a priority split rather than an
exclusion: the records it does not select remain in the frame and the resulting uncertainty
is discussed in Section~\ref{sec:threats}.

\subsection{Screening}
\label{sec:method:screening}

The frozen protocol prescribes two outcome-blinded human reviewers with a third resolving
criterion-level disagreement. Human coding was excluded for this project, so we used two LLM
model axes. For the 58-record blinded comparison, what is preserved is independence
between axes, blinding to the outcome, and disagreement surviving as \textsc{unresolved}
rather than being averaged into a verdict. What is not preserved is that the reviewers are
human, and that the second axis covers every record blindly. Both are recorded as deviations.

The screen separates scope from eligibility. Criteria
that turn on what a paper did -- whether its task identities are identifiable, whether its
protocol is reconstructable -- are not decidable from an abstract, and requiring them at
abstract level made the confirmatory verdict nearly unreachable. Peer-review status is a
metadata fact rather than a judgement and is decided from record type and venue, with an
explicit list of hosts that are institutional repositories rather than venues.

A preprint is not a dead end. The protocol deduplicates a preprint into its study family,
and one confirmed study is present in the pool only as its arXiv record because PMLR is
under-indexed by every channel that answered. Sibling resolution recovers the peer-reviewed
version by title from DBLP; treating an unresolved preprint as excluded would turn a
metadata gap into an eligibility decision.

The first axis judged all 929 priority records. The second, a different model family, judged
58 records blind: 38 first-axis \texttt{IN\_SCOPE} records and 20 decoys drawn from the first
axis's \texttt{OUT\_OF\_SCOPE} bucket, shuffled so the second axis could not tell them apart. After sibling
resolution, it judged 19 additional first-axis in-scope records in a non-blind coverage
batch. That batch extended coverage but is not independent axis-quality evidence. Agreement
in the blinded comparison was 45 of 58, with most disagreement being one axis uncertain
where the other was definite.

The resulting roster is 47 studies whose records were judged in scope by both model axes,
with a peer-reviewed record established from metadata or recovered by sibling resolution.
The independence and blinding evidence is limited to the 58-record blinded batch. It is a bounded
discovery cohort, not a census, and nothing in this paper treats it as one. Studies are
counted once at the strongest evidence band they earn, because one study appears as several
records and deduplicating within bands alone would count it twice.

\subsection{Where the four deep-audit studies came from}
\label{sec:method:provenance}

The four studies audited in depth were not sampled. They were inherited from an earlier
project, under a selection rule that is not documented and cannot be reconstructed, and this
paper's discovery protocol was written afterwards and around them: the roster of 47 is a
re-derivation of the frame those four came from, not an independent frame that happens to
contain them. What can be established is order, not independence -- the four were fixed
before any result reported here was produced -- and order alone does not rule out that the
earlier project chose them for artifact richness or for known reproduction difficulty. We
therefore claim no outcome-independent selection and no representativeness, and every claim
about the four is stated about those four.

\begin{table}[t]
\caption{What was frozen, when, and with what digest. A row marked \emph{not frozen} is
recorded as such rather than described as fixed.}
\label{tab:freeze}
\small
\begin{tabular}{p{0.23\columnwidth}p{0.22\columnwidth}p{0.38\columnwidth}}
\toprule
Object & Frozen & Evidence \\
\midrule
Protocol document & before item derivation & sha256 recorded as the derivation's input \\
Required-item vector (206) & before any availability lookup & status recorded in the vector file \\
Roster of 47 & before outcome coding & sha256 recorded in the register \\
Two-tier split & before any tier-1 result & recorded in the register \\
Four-study selection rule & \emph{not frozen} & inherited; rule not recoverable \\
\bottomrule
\end{tabular}
\end{table}

The 206 requirements were instantiated by applying the codebook to the four studies during
case application, and the vector carries the status \textsc{frozen before any availability
lookup} together with the protocol digest it was derived from. Table~\ref{tab:freeze} lists what was frozen, when, and against which digest. We disclose the frozen vector,
its protocol hash, and all retained provenance; per-requirement derivation links, timestamps
and revision histories were not retained, so we claim set-level priority and not an
item-level derivation trail.

\subsection{Two-tier audit}
\label{sec:method:tiers}

Auditing 47 studies at the depth applied to four is roughly ten times the work, and the cost
is not retrieval but the per-study closure judgement and axis derivation. The audit
therefore splits, and the split was frozen before any tier-1 result was inspected.

Tier one is the availability design across all 47: enumerate the artifact locators each paper itself
exposes, attempt them, and judge whether the deposit is the paper's own, whether it holds
run outputs, and whether anything in it states which released artifact produced which
reported number. Retrieval is automated and the three judgements are read, from a blind
evidence dump per study in which every automated verdict is withheld
(\S\ref{sec:method:verification}). Tier two was specified as compatibility depth on a sample
drawn after tier one and within tier-one availability classes at
$\min(4, |\text{class}|)$ per class. The draw used seed 20260809 and selected 18 non-census
studies; together with the four inherited census-stratum studies, the tier-two set contains
22 studies. No per-class compatibility analysis was run. The completed
corpus-scale reporting is the availability/disposition screen, including the 22 rostered
papers that expose no locator.

\subsection{Why the corpus claim is read rather than pattern-matched}
\label{sec:method:verification}

An automated screen was built for availability evidence first and checked against the four
inherited studies whose answers had been established by reading. A separate availability-
instrument check used channels the screen does not read. Neither check was a tier-two
compatibility analysis. Both changed the screen, and a third measurement retired it as the
instrument.

Validation caught the screen counting a source file and a test fixture as run outputs, and
reading an instruction to regenerate a figure as a binding to a reported number. The second
was not removed. Regenerating a figure from released data is weaker than stating which run
produced which number, but it is real, and deleting it to protect the hypothesis would have
been the wrong move; it is reported as its own grade.

The separate availability-instrument check found two error classes that together would have produced a
cleaner result than the truth. Locators harvested from full text included the benchmark
repositories a paper evaluates on rather than the artifact it published, so for one study
the screen audited three unrelated repositories while the study's own never entered the
list. And the screen read only the default branch, so artifacts attached to releases were
invisible -- one study attaches a 1.2\,GB run archive to a release and keeps nothing in the
tree. Both were corrected and the corpus rescreened.

The corrected availability screen was then measured against blind readings of the same evidence on 13
studies held out from every pattern fix, and it disagreed on 4 of 13 -- every binding disagreement in
the direction of missing one. Because the studies available to tune on are the studies the
claim is about, tightening further would have moved the error rate without measuring it. All
25 studies exposing a locator were therefore judged by reading for the availability screen,
each from a dump of its deposit's listing and prose with the screen's verdicts withheld, and the screen's verdicts
recorded beside the judgement afterwards; it disagrees on 8 of 25.
Section~\ref{sec:results:corpus} reports the read figures,
Section~\ref{sec:threats} what the screen got wrong, and the earlier screen-derived figures
are withdrawn.
